\documentclass[12pt]{article}

\usepackage[a4paper, margin=1in]{geometry}
\usepackage{amsmath, amssymb}
\usepackage{setspace}
\usepackage{hyperref}

\setstretch{1.2}

\title{\textbf{BrainGNN: \\ A Plastic Graph Neural Network Architecture for Node Classification}}
\author{Iyer Ramkumar Ramasubramanian \\ \textit{Independent Researcher}}
\date{\today}

\begin{document}

\maketitle

\begin{abstract}
Static neural network weights often fail to capture the dynamic adaptability of biological neural systems. This paper introduces \textbf{BrainGNN}, a Graph Neural Network architecture integrated into the \textbf{ULTRA Research Platform}. BrainGNN incorporates learned plasticity mechanisms, allowing network weights to modulate dynamically during the training process. Evaluated across standard benchmark datasets including Cora, CiteSeer, and PubMed, the model demonstrates how bio-inspired plasticity can enhance representational capacity in graph-structured data.
\end{abstract}

\section{Introduction}
As an independent researcher investigating the convergence of neurobiology and artificial intelligence, I propose that learning should be viewed as a continuous process of structural and weight adaptation. The BrainGNN model extends the principles of the \textit{Homo Deus Brain Model} by moving from reactive equilibrium to active, plastic learning within graph-based environments.

\section{System Architecture}

The BrainGNN framework utilizes a tiered approach to process graph-structured information:

\begin{itemize}
\item \textbf{Plastic Cognitive Layer}: A linear or graph-based layer (GCN/GAT) that implements a learned \texttt{plastic\_gate} parameter to modulate weight updates.
\item \textbf{Multi-Head Attention (GAT)}: Utilizes attention mechanisms to allow nodes to selectively weight the influence of their neighbors, mirroring synaptic entanglement.
\item \textbf{Automated Research Pipeline}: Includes distributed sweeps and automated \LaTeX\ paper generation for rapid experimental iteration.
\end{itemize}

\section{Mathematical Formulation: Learned Plasticity}

Unlike traditional networks where weights $W$ are updated solely via gradient descent, BrainGNN introduces a stochastic plasticity term modulated by a learned gate:

\[
P = \sigma(\text{plastic\_gate})
\]
\[
W_{t+1} = W_t + \eta \nabla L + \epsilon \cdot P
\]

Where $\epsilon$ represents a noise tensor and $\sigma$ is the sigmoid activation function. This mechanism allows the network to explore the weight space more dynamically, potentially avoiding local minima.

\section{Experimental Configuration}

The model is trained using the following hyperparameters across the Cora, CiteSeer, and PubMed datasets:

\begin{itemize}
\item \textbf{Optimizer}: Adam with a learning rate of 0.01.
\item \textbf{Architecture}: 32 hidden dimensions with 50 training epochs per run.
\item \textbf{Plasticity Strength}: Initialized with a coefficient of 0.005.
\end{itemize}

\section{Results and Inference}

Initial results generated by the \textbf{ULTRA Platform} indicate:
\begin{itemize}
\item \textbf{Accuracy}: The model achieves competitive classification accuracy on the test masks of standard citation networks.
\item \textbf{Dynamic Adaptation}: The inclusion of the plasticity gate enables the system to maintain representational diversity throughout the training cycle.
\item \textbf{Stability}: Despite the stochastic nature of the plasticity term, the model converges reliably across multiple runs.
\end{itemize}

\section{Conclusion}
BrainGNN demonstrates that integrating biological concepts like plasticity into GNN architectures can yield robust learning systems. By utilizing the ULTRA platform's automated capabilities, this research provides a scalable path for investigating high-dimensional, adaptive neural structures.

\section{Repository for Experimentation}
\noindent
\url{https://github.com/spacetracker-collab/Braingnn}

\end{document}
